\documentclass[11pt]{article}
\usepackage[a4paper,margin=21mm]{geometry}
\usepackage{fontspec}
\setmainfont{Latin Modern Roman}
\newfontfamily\CJKfont{Noto Serif CJK SC}
\XeTeXlinebreaklocale "zh"
\XeTeXlinebreakskip=0pt plus 1pt
\usepackage{amsmath,amssymb,amsthm,amscd,graphicx,color}
\usepackage{xurl}
\usepackage[unicode,hidelinks]{hyperref}
\allowdisplaybreaks
\setlength\emergencystretch{3em}
\numberwithin{equation}{section}

\usepackage{mathtools}

\numberwithin{equation}{section}

\newtheorem{Theorem}{Theorem}[section]
\newtheorem*{Theorem*}{Theorem}
\newtheorem{Corollary}[Theorem]{Corollary}
\newtheorem{Lemma}[Theorem]{Lemma}
\newtheorem{Proposition}[Theorem]{Proposition}
 { \theoremstyle{definition}
\newtheorem{Definition}[Theorem]{Definition}
\newtheorem{Note}[Theorem]{Note}
\newtheorem{Example}[Theorem]{Example}
\newtheorem{Remark}[Theorem]{Remark}
\newtheorem{Question}[Theorem]{Question} }

\begin{document}
\begin{center}\Large Countable mixed-identity-free groups embed in maximal discrete subgroups of their projective group-von-Neumann unitary groups in the trace topology\end{center}
\noindent\textbf{Stable ID:} 1248\quad\textbf{Author:} Vadim Alekseev; Andreas Thom\par
\noindent\textbf{Work:} Maximal Discrete Subgroups in Unitary Groups of Operator Algebras\par
\noindent\textbf{Source:} \url{https://sigma-journal.com/2022/052/}\par
\noindent\textbf{Version:} Official SIGMA 18 (2022) article 052, DOI 10.3842/SIGMA.2022.052; journal PDF and native TeX acquired 2026-09-30, exact bytes pinned by SHA256\par
\noindent\textbf{Rights:} CC BY 4.0. Authors retain copyright. \url{https://creativecommons.org/licenses/by/4.0/}. Reuse and typesetting adaptation; original language and mathematical content preserved.\par
\noindent\textbf{Difficulty (prior selection preserved):} {\CJKfont H2: The proof first derives two-sided free commutator trace estimates and iterates them while retaining a positive projective lower bound. It transfers asymptotic group freeness to a tracial ultrapower, then back to nontrivial words arbitrarily near the identity, establishing a uniform discrete radius for every overgroup. That common radius is the essential new step allowing the final chain-union maximality argument in a non-locally compact unitary group.}\par
\medskip\noindent\textbf{Editorial boundary note (not author text):} Selected complete author Theorem3.5 for countable mixed-identity-free groups, with the original full printed statement retained. The topology is the projective trace2-norm metric, not the later C*-operator-norm topology. Full trace expansion/lower-upper commutator bounds, iterated free-Haar word estimates, precise MIF/asymptotic-freeness prerequisite, complete author ultrapower enumeration argument, approximation proof, uniform-discreteness consequence and chain/Zorn proof are unchanged. Countability is the previously reviewed selected scope matching the explicit enumeration. Auxiliary free estimates and hyperfinite/C*-variants receive no separate counts. Every original literal inequality/index/wording is preserved without mathematical refereeing or editor correction.\par
\medskip\hrule\medskip\noindent\textbf{Author text begins below.}\par

% Original source lines 70--70
We aim to study this phenomenon in the context of finite von Neumann algebras, where the unitary group is metrized by the $2$-norm and no contractivity of the commutator map is available in general. We will make use of free probability theory to overcome this problem in the situation of the group von Neumann algebra $LG$, and show in many cases that nevertheless the class of discrete subgroups of the projective unitary group ${\rm PU}(LG)$ containing $G$ is uniformly discrete. As a consequence, $G$ is contained in a maximal discrete subgroup. We do not know if or when $G \subset {\rm PU}(LG)$ is maximal discrete itself.

\setcounter{section}{1}
% Original source lines 72--145
\section{Preparations using free probability}

Let $(M,\tau)$ be a finite von Neumann algebra and we will typically assume that it is a ${\rm II}_1$-factor. The $2$-norm on $M$ is defined as $\|x\|_2 = \tau(x^*x)^{1/2}$. We denote its unitary group by ${\rm U}(M)$ and its projective unitary group by ${\rm PU}(M):={\rm U}(M)/S^1$. Consider unitaries $u,v \in {\rm U}(M)$ which are freely independent, that is to say that they lie in freely independent subalgebras in the sense of Voiculescu \cite{MR1217253}. We set $\alpha := \tau(u)$ and $\beta := \tau(v)$. Note that because of freeness we have
\begin{gather*}
0 = \tau((u-\alpha)(v-\beta)) = \tau(uv) - \beta \tau(u) - \alpha \tau(v) + \alpha \beta
\end{gather*}
and hence $\tau(uv) = \tau(u)\tau(v)$ for any freely independent pair of unitaries $u,v \in {\rm U}(M)$. We aim to compute the trace of $uvu^*v^*$. Again, by freeness, we obtain
\begin{align*}
0 ={}& \tau((u-\alpha)(v-\beta)(u^* - \bar \alpha)(v^* - \bar \beta))
\\
={}& \tau(uvu^*v^*) - \alpha \tau(vu^*v^*) - \beta \tau(uu^*v^*) - \bar \alpha \tau(uvv^*) - \bar \beta \tau(uvu^*)
\\
& + \alpha \beta \tau(u^*v^*) + |\alpha|^2 \tau(vv^*) + \alpha \bar \beta \tau(vu^*) + \bar \alpha \beta \tau(uv^*) + |\beta|^2 \tau(uu^*) + \bar \alpha \bar \beta \tau(uv)
\\
& - |\alpha|^2 \beta \tau(v^*) - \alpha |\beta|^2 \tau(u^*) - |\alpha|^2 \bar \beta \tau(v) - \bar \alpha |\beta|^2 \tau(u) + |\alpha|^2|\beta|^2
\\
={}& \tau(uvu^*v^*) - 2|\alpha|^2 - 2 |\beta|^2 + 4 |\alpha|^2 |\beta|^2 + |\alpha|^2 + |\beta|^2 - 4 |\alpha|^2 |\beta|^2 + |\alpha|^2 |\beta|^2
\\
={}& \tau(uvu^*v^*) - |\alpha|^2 - |\beta|^2 + |\alpha|^2 |\beta|^2.
\end{align*}
Hence, we conclude that
\[
\tau(uvu^*v^*) = |\alpha|^2 + |\beta|^2 - |\alpha|^2 |\beta|^2 = 1 - \big(1-|\tau(u)|^2\big)\big(1-|\tau(v)|^2\big)
\] for any freely independent pair of unitaries $u,v \in {\rm U}(M)$.
On ${\rm U}(M)$, we consider the length function $\ell(u):= \|1-u\|_2 = \sqrt{2-2 \operatorname{Re}(\tau(u))}$.
Since
\begin{gather*}
(1-|\tau(u)|^2) = (1+ |\tau(u)|)(1-|\tau(u)|) \leq 2 (1-|\tau(u)|) \leq 2 - 2\operatorname{Re}(\tau(u)) = \ell(u)^2,
\end{gather*}
we obtain
\begin{align*}
\ell([u,v]) &= \sqrt{2-2 \operatorname{Re}(\tau(uvu^*v^*))} \\
&= \sqrt{2\big(1-|\tau(u)|^2\big)\big(1- |\tau(v)|^2\big)} \\
&\leq\sqrt{2} \cdot \ell(u) \ell(v)
\end{align*}
for any freely independent pair of unitaries $u,v \in {\rm U}(M)$. With only little more work, we also obtain a lower bound for $\ell([u,v])$.
Indeed, define the projective length function
\begin{gather*}
\bar \ell(u) = \inf_{|\lambda|=1 } \|\lambda - u\|_2 = \sqrt{2(1-|\tau(u)|)},
\end{gather*}
which metrizes ${\rm PU}(M)$ in a natural way. Note that
$\mu - \mu^2/4 \geq 1/2 \cdot \mu$ for $\mu \in [0,2]$. If $u$, $v$ are freely independent unitaries, we get
\begin{align*}
\ell([u,v])^2 &= 2(1 - \operatorname{Re} \tau([u,v])) \\
&= 2 \big(1-|\tau(u)|^2\big)\big(1- |\tau(v)|^2\big) \\
&= 2 \big(1-\big(1-\bar\ell(u)^2/2\big)^2\big) \big(1-\big(1-\bar\ell(v)^2/2\big)^2\big) \\
&= 2 \big( \bar \ell(u)^2 - \bar \ell(u)^4/4\big) \big( \bar \ell(v)^2 - \bar \ell(v)^4/4\big) \\
&\geq 2 (1/2)^2 \cdot \ell(u)^2 \ell(v)^2.
\end{align*}
Hence, we conclude that
\begin{gather*}
\ell([u,v]) \geq 1/\sqrt{2} \cdot \bar \ell(u) \bar \ell(v)
\end{gather*}
and we may summarize the computations as follows:
\begin{Theorem}
Let $u,v \in {\rm U}(M)$ be freely independent unitaries. Then, we have
\begin{gather*}
1/\sqrt{2} \cdot \bar \ell(u) \bar \ell(v) \leq \bar\ell([u,v])= \ell([u,v]) \leq \sqrt{2} \cdot \bar \ell(u) \bar \ell(v) \leq \sqrt{2} \cdot \ell(u) \ell(v).
\end{gather*}
\end{Theorem}

For the rest of the paper we fix a sequence of words $(w_n)_n$ in $F_2= \langle x,y \rangle$ such that $w_1=x$ and inductively $w_{n+1} = [w_n,y^{n}xy^{-n}]$, i.e., $w_2=\big[x,yxy^{-1}\big]$, $w_3=\big[\big[x,yxy^{-1}\big],y^2xy^{-2}\big],$ and so~on. If $u$, $v$ are freely independent and $v$ is a Haar unitary (i.e., a unitary with $\tau(v^n)=0$ for all $n \in \mathbb Z \setminus \{0\}$), it is easy to see that $ u,vuv^*,v^2u(v^*)^2,\dots$ are freely independent too. We consider the sequence $(w_n(u,v))_n$, for $i \in \mathbb N$, i.e.,
\begin{gather*}
u_1=u, \qquad
u_2=[u,vuv^*], \qquad
u_3=\big[[u,vuv^*],v^2u(v^*)^2 \big],\qquad \dots\,.
\end{gather*}
Thus, we obtain from the previous theorem by induction
\begin{gather*}
\big(1/\sqrt{2}\big)^{n-1} \cdot \bar\ell(u)^n \leq \ell(w_n(u,v)) \leq \big(\sqrt{2}\big)^{n-1} \cdot \ell(u)^n.
\end{gather*}

The upper bound was rather surprising for us at first, since $\ell(u)< 1/\sqrt{2}$ does not seem to be a severe assumption and we may assume in addition that the subgroup generated by $u$ in ${\rm PU}(M)$ is discrete, we may even assume that $u'$ is a Haar unitary in a corner $pMp$ and $u = u' + p^{\perp}$ for a projection $p$ with $\tau(p)<1/4$.



% Original source lines 163--226
\section{Applications to group von Neumann algebras}

We will now turn our attention to group von Neumann algebras. Let $G$ be a group. A sequ\-ence~$(g_n)_n$ of elements in $G$ is called asymptotically free if $(g_n)_n$ is free from the diagonal copy of $G$ in $\prod_n G$. Equivalently, for every $k \in \mathbb N$, $s_1,\dots,s_k \in G \setminus \{1\}$ and every $\varepsilon_1,\dots,\varepsilon_k \in \mathbb Z \setminus \{0\}$, there exists $n \in \mathbb N$, such that
\begin{gather*}
s_1 g_n^{\varepsilon_1} s_2 g_n^{\varepsilon_2} s_3 \cdots g_n^{\varepsilon_k} \neq 1.
\end{gather*}
Groups admitting an asymptotically free sequence have been studied for a long time.

Recall that a mixed identity for $G$ is a word $w \in \mathbb Z \ast G = \langle G, t\rangle$ in one variable $t$ with coefficients in $G$ such that $w$ evaluates to the identity if any element of $G$ is substituted for $t$. A~group is called mixed identity free (MIF) if it does not satisfy any non-trivial mixed identity. We~refer the reader to \cite{MR3556970} for further details around this notion, emphasizing the following statement.

\begin{Proposition}[{\cite[proof of Proposition 5.3 and Remark 5.1]{MR3556970}}]
Let $G$ be a group. The following are equivalent:
\begin{enumerate}\itemsep=0pt
\item[$1.$] The group $G$ contains an asymptotically free sequence.
\item[$2.$] There is no mixed identity $w \in \mathbb Z \ast G$ for $G$.
\item[$3.$] $G$ is mixed identity free, i.e., there exists no mixed identity in $\mathbb F_\infty \ast G$.
\end{enumerate}
\end{Proposition}

We denote the group von Neumann algebra of $G$ with its natural trace by $(LG,\tau)$. It is well-known that any non-trivial conjugacy class in a ${\rm MIF}$~group is infinite. In particular, the group von Neumann algebra $LG$ of a ${\rm MIF}$~group $G$ is a ${\rm II}_1$-factor. Note that the left-regular representation provides a natural inclusion $G \subset {\rm PU}(LG)$.
 We will be interested in the existence if maximal discrete subgroups of ${\rm PU}(LG)$ containing $G$.

\begin{Lemma}
Let $(g_n)_n$ be an asymptotically free sequence. There exists an ultrafilter $\omega$, such that $[(g_n)_n] \in U((LG)^{\omega})$ is freely independent from the diagonal copy of $LG$ in $(LG)^\omega$.
\end{Lemma}
This lemma is standard but not completely trivial, and we include the argument for the sake of completeness.
\begin{proof}
Let $(w_m)$ be an enumeration of all mixed identities, that is, elements in $\mathbb Z\ast G$. We~denote by $w(m_1,\dots,m_\ell)$ the iterated commutator
\begin{gather*}
w(m_1,\dots,m_\ell) = [w_{m_1},[w_{m_2},\dots, [w_{m_{\ell-1}},w_{m_\ell}]]].
\end{gather*}
Observe that if $w(m_1,\dots,m_\ell)(g) \neq 1$, it guarantees that $w_{m_i}(g)\neq 1$ for all $i=1,\dots,\ell$. Let $v_\ell\coloneqq w(1,2,\dots,\ell)$. Asymptotic freeness of $(g_n)$ guarantees the existence of natural numbers $n_\ell$ such that $v_\ell(g_{n_\ell})\neq 1$. Taking an ultrafilter containing this sequence finishes the proof.
\end{proof}

\begin{Lemma}
Let $G$ be a ${\rm MIF}$~group and $u \in {\rm U}(LG)$ with $u \not \in S^1 \cdot 1$ and $\operatorname{Re}(\tau(u))>3/4$. For any $\varepsilon>0$ there exists some $w \in \mathbb F_2$ and $g \in G$ such that $\ell(w(u,g))< \varepsilon$. Moreover, the image of the map from $\mathbb Z \ast G$ to ${\rm PU}(LG)$ that maps the generator $1 \in \mathbb Z$ to $u$ is not discrete.
\end{Lemma}
\begin{proof}
Note that from the assumptions $\bar \ell(u) \neq 0$ and $\ell(u)< \sqrt{2}$.
Let $(g_k)_k$ be an asymptotically free sequence and set $v= [(g_k)_k] \in (LG)^{\omega}$. By assumption $\ell(u) \leq 1/\sqrt{2} - \delta$ for some $\delta>0$. By~the results from above, we get
\begin{gather*}
\big(1/\sqrt{2}\big)^{n-1} \cdot \bar\ell(u)^n \leq \ell(w_n(u,v)) \leq\big(\sqrt{2}\big)^{n-1} \cdot \ell(u)^n.
\end{gather*}
Thus, for some $k \in \mathbb N$
\begin{gather*}
\big(1/\sqrt{2}\big)^{n} \cdot \bar\ell(u)^n \leq \ell(w_n(u,g_k)) \leq \big(\sqrt{2}\big)^{n} \cdot \ell(u)^n
\end{gather*}
and this finishes the proof.
\end{proof}

The following consequence is immediate.
\begin{Corollary}
Let $G$ be a ${\rm MIF}$~group and let $u \in {\rm PU}(LG)$ with $0 < \bar \ell(u) < 1/\sqrt{2}$. Then, the subgroup $\langle u, G \rangle \subset {\rm PU}(LG)$ is not discrete.
In particular, the set of discrete subgroups of~${\rm PU}(LG)$ containing $G$ is uniformly discrete.
\end{Corollary}
\begin{proof} By the previous lemma, no such subgroup contains a non-trivial element with $\bar \ell(u)<1/\sqrt{2}$. Indeed, we would apply the previous lemma to a suitable lift $u'$ with $\ell(u')< 1/\sqrt{2}$.\end{proof}

We are now able to state and prove our main result.
\begin{Theorem} %\label{main}
Let $G$ be a ${\rm MIF}$~group. Then, $G$ is contained in a maximal discrete subgroup in~${\rm PU}(LG)$.
\end{Theorem}
\begin{proof}
The discrete subgroups of ${\rm PU}(LG)$ containing $G$ are ordered by inclusion. Now, the result follows by Zorn's lemma since unions over chains remain discrete.
\end{proof}
\begin{thebibliography}{99}
\bibitem[6]{MR3556970}
Hull M., Osin D., Transitivity degrees of countable groups and acylindrical
 hyperbolicity, \href{https://doi.org/10.1007/s11856-016-1411-9}{\textit{Israel~J. Math.}} \textbf{216} (2016), 307--353,
 \href{https://arxiv.org/abs/1501.04182}{arXiv:1501.04182}.

\bibitem[12]{MR1217253}
Voiculescu D.V., Dykema K.J., Nica A., Free random variables, \textit{CRM
 Monograph Series}, Vol.~1, \href{https://doi.org/10.1090/crmm/001}{Amer. Math. Soc.}, Providence, RI, 1992.

\end{thebibliography}
\end{document}
